%% ============================================================
%%  Abstract & Keywords — sn-jnl / Nature Computational Science
%%  Target length: ≤150 words (NCS strict guideline)
%%  Style: "Here we show..." declarative opening
%% ============================================================

\abstract{%
As artificial intelligence models scale beyond fast-memory capacity,
inference latency is governed by coordination across hierarchical memory
systems rather than computation alone.
Here we show that this coordination is not continuously optimizable but
exhibits intrinsic, regime-dependent limits whose transitions are abrupt.
Expressing system configuration as two dimensionless ratios---the
fast-memory residency ratio ($R_C$) and the cross-tier transfer-pressure
ratio ($R_B$)---we identify three operationally distinct regimes:
capacity-limited, coordination-dominated, and I/O-limited.
Within coordination-dominated regimes, orchestration reduces end-to-end
latency by redistributing costs across computation, locality, transfer,
and synchronization.
Beyond capacity or bandwidth thresholds, the same strategies fail
predictably, and additional coordination becomes additive overhead.
We derive closed-form regime boundaries from a minimal multi-tier transfer
model and validate phase-like transitions across language and vision
workloads on GPU, TPU, and Inferentia-based systems.
These results reframe hierarchical memory optimisation as regime-aware
design rather than continuous tuning, with implications for any
memory-bound computing system.
}

\keywords{hierarchical memory systems, memory orchestration, AI inference,
phase transitions, regime-dependent behaviour, large language models,
computational systems}
